\section{Introduction}
\label{sec:intro}

Nocturnal stable boundary layers (SBLs) remain among the most poorly understood and inconsistently modeled features of the lower atmosphere \citep{Mahrt2014, Acevedo2016}. Unlike daytime convective boundary layers, where surface heating drives turbulence vertically, nocturnal SBLs exhibit a fundamental instability: strong temperature inversions suppress vertical mixing, yet persistent wind shear continues to generate turbulent kinetic energy (TKE) locally. The result is an environment dominated by \emph{intermittency}---alternating periods of quiescent, laminar flow interrupted by sudden shear-driven turbulent bursts that collapse and dissipate over minutes to tens of minutes \citep{Mahrt1999, Terradellas2001}.

From a statistical perspective, these profile trajectories are consistent with a Brownian-yet-non-Gaussian regime: variance can grow diffusively while the instantaneous profile-increment distribution remains heavy-tailed because the effective diffusivity is itself state dependent \citep{Sposini2018}. This motivates interpreting broad variance envelopes as signatures of intermittently switching transport states rather than as evidence of a single stationary Gaussian process.

Superimposed on this intermittency are horizontally propagating internal gravity waves (IGWs), generated by drainage flows over complex terrain or by baroclinic vorticity forcing at mesoscale boundaries \citep{Vosper2006, Sun2015}. As synthesized by \citet{Sun2015}, stably stratified environments naturally support an overlapping continuum of anisotropic, intermittent turbulence and wavelike motions. These waves transport energy horizontally and can locally steepen density gradients, triggering instability and turbulent collapse. However, traditional single-column decomposition schemes (e.g., Richardson number thresholds, wavelet filtering) treat waves and turbulence as completely separable, missing the \emph{coupled} dynamics that govern real boundary-layer transitions.

\begin{figure}[htbp]
    \centering
    % NOTE: If the figure does not rotate correctly on compile,
    % consider using the rotating package with \begin{sidewaysfigure}
    \includegraphics[width=0.75\linewidth, angle=-90]{Night_Sky_Turbulence_Classification_Framework.png.pdf}
    \caption{Classification framework for night sky stable boundary layer regimes mapping scale-dependent multi-regime transitions.}
    \label{fig:framework}
\end{figure}

\subsection{The MOST Crisis}

The foundational framework for stable boundary-layer parameterization remains Monin--Obukhov similarity theory (MOST) \citep{Obukhov1946, Paulson1970}, which predicts that normalized gradients and fluxes depend only on a single stability parameter:
\begin{equation}
\zeta = \frac{z}{L} = \frac{z \kappa g}{\theta_{\mathrm{ref}}} \frac{\overline{T'w'}}{u_*^2 \theta_{\mathrm{ref}}}
\label{eq:zeta}
\end{equation}
where $L$ is the Obukhov length scale, $u_*$ is friction velocity, and $\overline{T'w'}$ is the sensible heat flux. For stable conditions ($\zeta > 0$), MOST predicts:
\begin{equation}
\begin{aligned}
\phi_m(\zeta) &= 1 + \beta_m\zeta, \\
\phi_h(\zeta) &= 1 + \beta_h\zeta,
\end{aligned}
\end{equation}
with universal constants $\beta_m \approx 5$ and $\beta_h \approx 5$ (though empirical estimates vary widely across regimes, spanning $\beta_m \in [3, 10]$ and $\beta_h \in [3, 15]$ \citep{Dyer1974, Högström1988}).

However, observational datasets from tower campaigns---such as CASES-99 \citep{Poulos2002}, SHEBA \citep{Uttal2002}, and FINN \citep{Mahrt2013}---reveal systematic deviations from MOST scaling laws. Flux-profile relationships show distinct hysteresis loops: the same gradient Richardson number yields entirely different flux estimates depending on whether the system is transitioning into or out of an inversion layer. Furthermore, wave-affected periods show scaling parameter values far outside the universal range, with non-monotonic behavior in $\zeta$, while intermittent turbulence produces rapid switches in effective $L$ that render time-averaged diagnostics fundamentally ambiguous.

As established in the comprehensive review by \citet{Sun2015}, state-of-the-art numerical weather prediction models fail systematically to simulate this seemingly random turbulence intensification because they lack a rigorous mechanism to parameterize wave-turbulence interactions within stably stratified geophysical flows.

Accordingly, we formulate the estimator around structural events instead of continuous local gradients: in non-stationary boundary-layer records, event segmentation has stronger physical interpretability than gradient-only tracking when transitions are abrupt and scale-coupled \citep{Kang2014a, Narasimha1995}.

\subsection{Spectral Decomposition Approaches}

Recent efforts have attempted to isolate waves from turbulence using wavelet filtering \citep{Vosper2006, Sun2015} or empirical mode decomposition (EMD) \citep{Huang1998}. While wavelet filtering decomposes signals into time-frequency atoms, the basis choice (e.g., Morlet, Daubechies) remains fundamentally arbitrary and introduces non-local Gibbs ringing near physical discontinuities. EMD extracts intrinsic mode functions adaptively, but its outputs are non-orthogonal, and the decomposition lacks mathematical uniqueness. Similarly, Proper Orthogonal Decomposition (POD) and Dynamic Mode Decomposition (DMD) techniques \citep{Schmid2010} extract dominant modes cleanly but require dense spatiotemporal snapshot matrices, which are entirely unavailable from sparse vertical instrument masts.

None of these methods preserve the physical constraint that projections must respect the \emph{natural metric} of the atmospheric system. A wave acting on a stable density gradient operates under a different effective mass than a turbulent eddy in the same layer. Standard spectral projections treat both fields identically, obscuring the primary coupling modes highlighted by \citet{Sun2015}.

\subsection{Contribution: Metric-Consistent Manifold Decomposition}

This paper presents a fundamentally different approach: instead of projecting data onto arbitrary basis functions, we construct a \emph{Riemannian metric} $\mathbf{M}$ that weights projections according to physical density and stratification. The metric ensures that wave energy is cleanly localized to low-frequency Chebyshev modes that respect the stable density gradient, while turbulent kinetic energy is concentrated in high-frequency modes highly sensitive to localized vertical wind shear. 

Key innovations presented herein include:
\begin{enumerate}
    \item \textbf{Hyperbolic compactification} of the vertical coordinate: Non-uniform grid packing ($\Delta z_{\min} = \SI{2.85}{\milli\meter}$ at the surface, $\Delta z_{\max} = \SI{9.43}{\meter}$ aloft) captures acute \SI{5}{\kelvin\per\meter} inversions without wasting degrees of freedom on upper air layers.
    \item \textbf{Thin SVD-truncated matrix projection}: With 8 instrument levels but an $N=32$ order spectral expansion, rank deficiency is structurally unavoidable. We solve this inverse problem via an economy-size, thin Singular Value Decomposition (SVD), restricting matrix operations entirely to the data-constrained subspace.
    \item \textbf{Partition-of-unity spectral masking} $(\psi_M, \psi_W, \psi_T)$: This framework explicitly separates mesoscale mean flow, gravity-wave scales, and turbulent dissipation without introducing ad-hoc cutoff boundaries.
\end{enumerate}

We frame this study as a reduced-order representation of SBL state evolution motivated by observed time-scale separation between fast turbulence adjustment and slow mesoscale forcing. Within that scope, we link the diagnostic metric set (including $D_{\mathrm{eff}}$ as an empirical modal-occupancy proxy) to the Ozmidov buoyancy scale, provide a generalized ``Virtual Tower'' interpolation framework to resolve the grid-topology mismatch for global model validation, and supply a spectral anti-aliasing scheme necessary for forward-predictive closure experiments.